% !PRINTED 371
\weberpage{371}
Daraus folgt, mit Rücksicht auf
\[
\left(\frac{b-\sqrt D}{2}\right)^2 = a_1 a_2 c + b\,\frac{b-\sqrt D}{2}:
\]
\weq{11}{
\mu\nu = a_1 a_2 (u_1 v_1 - c\,u_2 v_2) + \frac{b-\sqrt D}{2}\,(a_1 u_1 v_2 + a_2 u_2 v_1 + b\,u_2 v_2).
}
Macht man also die bilineare Substitution:
\weq{12}{
\begin{aligned}
t_1 &= u_1 v_1 - c\,u_2 v_2, \\
t_2 &= a_1 u_1 v_2 + a_2 u_2 v_1 + b\,u_2 v_2,
\end{aligned}
}
so ergibt sich
\weq{13}{\lambda = \mu\nu,}
und folglich
\weq{14}{T = UV.}
In der Forderung, daß $a_1$ und $a_2$ relativ prim sein sollen, liegt bisweilen eine gewisse Unbequemlichkeit, z.\,B. wenn es sich um die Komposition einer Form mit sich selbst handelt.

Lassen wir also die Voraussetzung fallen, daß $\mathfrak{a}$ zu $\mathfrak{b}$ und $a_1$ zu $a_2$ relativ prim seien, halten dagegen an der Annahme fest, daß
\weq{15}{b^2 - D = 4\,a_1 a_2 c}
durch $4 a_1 a_2$ teilbar sei, so gilt die Relation (11) noch, und es folgt, daß $\mathfrak{a}\mathfrak{b}$ der größte gemeinschaftliche Teiler von
\[
a_1 a_2, \quad a_1\,\frac{b-\sqrt D}{2}, \quad a_2\,\frac{b-\sqrt D}{2}, \quad b\,\frac{b-\sqrt D}{2}
\]
ist (Bd.\,II, §~160, 2).

Haben nun $a_1$, $a_2$, $b$ keinen gemeinschaftlichen Teiler, so kann man der Gleichung $k_1 a_1 + k_2 a_2 + kb = 1$ durch ganze rationale $k$ genügen, und demnach ist $\mathfrak{a}\mathfrak{b}$ auch der größte gemeinschaftliche Teiler von
\weq{16}{a_1 a_2, \quad \frac{b-\sqrt D}{2}.}

2. \gesperrt{Es ist zu beweisen, daß $a_1 a_2$ die kleinste durch $\mathfrak{a}\mathfrak{b}$ teilbare natürliche Zahl ist.}

Setzen wir nämlich diese kleinste Zahl $=a$, so ist $a_1 a_2$ durch $a$ teilbar, etwa
\[
a_1 a_2 = a\,a'.
\]
